\documentclass[../../../include/open-logic-section]{subfiles}

\begin{document}

\olfileid{sth}{cardinals}{classing}
\olsection{Winates, \usetoken{S}{enumerable}, \usetoken{S}{nonenumerable}}

Sawisé ngenal kardinal, prayoga kita ngrembug sawatara jinis
kardinal, mliginé kardinal winates, !!{enumerable}, lan
!!{nonenumerable}.

Rong asil kapisan bakal nuduhaké yèn kardinal winates persis
ordinal winates, kang wis kita tetapake minangka
\emph{wilangan asli} ing
\olref[z][infinity-again]{defnomega}:

\begin{prop}\ollabel{finitecardisoequal}
Upamané $n, m \in \omega$. Mula $n = m$ yen lan mung yen $\cardeq{n}{m}$.
\end{prop}

\begin{proof}
\emph{Saka kiwa menyang tengen} cetha. Kanggo mbuktèkaké
\emph{saka tengen menyang kiwa}, upamané $\cardeq{n}{m}$
nanging $n \neq m$. Miturut Trikotomi, $n
\in m$ utawa $m \in n$; tanpa ngurangi keumuman, upamané $n \in m$.
Mula $n \subsetneq m$ lan ana !!a{bijection} $f \colon m \to n$,
saéngga $m$ tanpa wates miturut Dedekind, mbantah
\olref[z][infinity-again]{naturalnumbersarentinfinite}.
\end{proof}

\begin{cor}\ollabel{naturalsarecardinals}
Yèn $n \in \omega$, mula $n$ iku kardinal.
\end{cor}

\begin{proof}
Langsung.
\end{proof}
\noindent
Mula sawetara pangertèn kang lumrah kanggo nyebut kardinal
“winates” utawa “tanpa wates” uga ekivalen:
\begin{thm}\ollabel{generalinfinitycharacter}Kanggo sembarang himpunan $A$, pratelan iki ekivalen:
\begin{enumerate}
	\item\ollabel{card:notinomega} $\card{A} \notin \omega$, yaiku
	$A$ dudu himpunan winates;
	\item\ollabel{card:omegaplus} $\omega \leq \card{A}$;
	\item\ollabel{card:infinite} $A$ tanpa wates miturut Dedekind.
\end{enumerate}
\end{thm}

\begin{proof}
Saka \olref[ord-arithmetic][using-addition]{ordinfinitycharacter},
\olref[cardinals][cardsasords]{lem:CardinalsBehaveRight}, lan
\olref{naturalsarecardinals}.
\end{proof}

Iki ngidinaké kita netepaké sacara \emph{resmi} sawetara
pangertèn kang sadurungé dienggo kanthi informal ing
\olref[sfr][][]{part}:

\begin{defn}\ollabel{defnfinite}
Kita kandha yèn $A$ \emph{winates} yen lan mung yen $\card{A}$
iku wilangan asli, yaiku $\card{A} \in \omega$. Yèn ora,
kita kandha yèn $A$ \emph{tanpa wates}.
\end{defn}
\noindent
Nanging wigatèkna yèn definisi iki gumantung marang
$\ZFC$. Kita perlu Pengurutan Becik kanggo njamin
saben himpunan nduwèni kardinalitas. Malah tanpa
Pengurutan Becik, bisa ana himpunan kang ora winates
nanging uga ora tanpa wates miturut Dedekind. Bab kaya iki
bakal kita rembug manèh ing \olref[choice][]{chap}.
Saiki kita terus nganggep Pengurutan Becik.

Saiki ayo ngalih saka kardinal winates menyang kardinal
tanpa wates. Iki rong pratelan dhasar:

\begin{cor}\ollabel{omegaisacardinal}
$\omega$ iku kardinal tanpa wates kang paling cilik.
\end{cor}

\begin{proof}
$\omega$ iku kardinal, amarga $\omega$ tanpa wates
miturut Dedekind, lan yèn $\cardeq{\omega}{n}$
kanggo sawatara $n \in \omega$, mula $n$ uga tanpa wates
miturut Dedekind, kang mbantah
\olref[z][infinity-again]{naturalnumbersarentinfinite}. Saiki
$\omega$ iku kardinal tanpa wates paling cilik miturut definisi.
\end{proof}

\begin{cor}
Saben kardinal tanpa wates iku ordinal limit.
\end{cor}

\begin{proof}
Upamané $\alpha$ ordinal penerus kang tanpa wates, mula
$\alpha = \beta
\ordplus 1$ kanggo sawatara $\beta$. Amarga penerus ordinal
winates mesthi winates, $\beta$
uga tanpa wates, mula $\cardeq{\beta}{\beta \ordplus 1}$ miturut
\olref[ord-arithmetic][using-addition]{ordinfinitycharacter}. Saiki
$\card{\beta} = \card{\beta\ordplus 1} = \card{\alpha}$ miturut
\olref[cardinals][cardsasords]{lem:CardinalsBehaveRight}; amarga
kardinal sawijining ordinal ora ngluwihi ordinal kasebut, mula
$\alpha \neq \card{\alpha}$.
\end{proof}

Ing \olref[sfr][siz][enm-alt]{defn:enumerable} kita wis nuduhaké
yèn himpunan tanpa wates kang !!{enumerable} bisa dibédakaké
saka kang !!{nonenumerable}. Definisi iku nuwuhaké asil iki:

\begin{prop}
$A$ iku !!{enumerable} yen lan mung yen $\card{A} \leq \omega$,
lan $A$ iku !!{nonenumerable} yen lan mung yen $\omega < \card{A}$.
\end{prop}

\begin{proof}
Miturut Trikotomi, rong pratelan mau ekivalen; mula cukup
mbuktèkaké yèn $A$ iku !!{enumerable} yen lan mung yen
$\card{A} \leq \omega$. Kanggo
\emph{saka tengen menyang kiwa}: yèn $\card{A} \leq \omega$, mula
$\cardle{A}{\omega}$ miturut
\olref[cardinals][cardsasords]{lem:CardinalsBehaveRight} lan
\olref{omegaisacardinal}. Kanggo \emph{saka kiwa menyang tengen}:
upamané $A$ iku !!{enumerable}; miturut
\olref[sfr][siz][enm-alt]{defn:enumerable} ana telung kasus:
\begin{enumerate}
	\item yèn $A = \emptyset$, mula $\card{A} = 0 \in \omega$, miturut
	\olref{naturalsarecardinals} lan
	\olref[cardinals][cardsasords]{lem:CardinalsBehaveRight}.
	\item yèn $\cardeq{n}{A}$, mula $\card{A} = n \in \omega$, miturut
	\olref{naturalsarecardinals} lan
	\olref[cardinals][cardsasords]{lem:CardinalsBehaveRight}.
	\item yèn $\cardeq{\omega}{A}$, mula $\card{A} = \omega$, miturut \olref{omegaisacardinal}.
\end{enumerate}
Dadi ing kabèh kasus, $\card{A} \leq \omega$.
\end{proof}
\noindent
Pancen, $\omega$ nduwèni kalungguhan mligi. Sanadyan ana
akèh ordinal kang bisa diétung:

\begin{cor}
$\omega$ iku siji-sijiné kardinal tanpa wates kang !!{enumerable}.
\end{cor}

\begin{proof}
Upamané $\cardfont{a}$ iku kardinal tanpa wates kang !!a{enumerable}. Amarga
$\cardfont{a}$ tanpa wates, $\omega \leq \cardfont{a}$. Amarga
$\cardfont{a}$ iku kardinal kang !!a{enumerable}, $\cardfont{a} =
\card{\cardfont{a}} \leq \omega$. Mula $\cardfont{a} = \omega$ miturut
Trikotomi. \end{proof}

Mesthiné ana tanpa wates akèhé kardinal. Mula kita bisa takon:
\emph{Pira cacahé kardinal?} Asil-asil sabanjuré nuduhaké
yèn pitakonan iki perlu dipikir manèh.

\begin{prop}\ollabel{unioncardinalscardinal}
Yèn saben unsur $X$ iku kardinal, mula $\bigcup X$ iku kardinal.
\end{prop}

\begin{proof}
Gampang dipriksa yèn $\bigcup X$ iku ordinal. Upamané $\alpha \in
\bigcup X$ iku ordinal; mula $\alpha \in \cardfont{b} \in X$ kanggo
sawatara kardinal $\cardfont{b}$. Amarga $\cardfont{b}$ iku kardinal,
$\cardless{\alpha}{\cardfont{b}}$. Amarga $\cardfont{b} \subseteq
\bigcup X$, kita nduwèni $\cardle{\cardfont{b}}{\bigcup X}$, lan mula
$\cardneq{\alpha}{\bigcup X}$. Kanthi nggeneralisasi, $\bigcup X$
iku kardinal.
\end{proof}

\begin{thm}\ollabel{lem:NoLargestCardinal}
Ora ana kardinal paling gedhé.
\end{thm}

\begin{proof}
Kanggo sembarang kardinal $\cardfont{a}$, Teorema Cantor
(\olref[sfr][siz][car]{thm:cantor}) lan
\olref[cardinals][cardsasords]{lem:CardinalsExist} ndadèkaké
$\cardfont{a} < \card{\Pow{\cardfont{a}}}$.
\end{proof}

\begin{thm}
Himpunan kabèh kardinal ora ana.
\end{thm}

\begin{proof}
Kanggo kontradiksi, upamané $C = \Setabs{\cardfont{a}}{\cardfont{a} \text{
iku kardinal}}$. Saiki $\bigcup C$ iku kardinal miturut
\olref{unioncardinalscardinal}, mula miturut \olref{lem:NoLargestCardinal}
ana kardinal $\cardfont{b} > \bigcup C$. Miturut definisi,
$\cardfont{b} \in C$, mula $\cardfont{b} \subseteq \bigcup{C}$, saéngga
$\cardfont{b} \leq \bigcup C$, sawijining kontradiksi.
\end{proof}

Bandhingna iki karo Paradoks Russell lan Burali-Forti.

\end{document}
